\section{Excluding one complete affine extension level}\label{sec:affine}

The product construction belongs to the extension-Nullstellensatz approach
of Buss, Impagliazzo, Kraj\'i\v{c}ek, Pudl\'ak, Razborov, and Sgall
\cite{BIKPRS}. We use a complete one-level family and prove its required
properties explicitly. An accuracy-\(h\) block has a finite tuple of old
affine polynomials \(g_1,\ldots,g_s\), fresh coefficient variables
\(r_{ui}\), and product
\begin{equation}\label{eq:block}
 P=\prod_{u=1}^h\left(1-\sum_{i=1}^s r_{ui}g_i\right).
\end{equation}
Its complete axioms are all companions \(g_iP\) and all coefficient
Boolean equations \(r_{ui}^2-r_{ui}\). The product \(P\) itself is not an
axiom. Coefficients of different blocks are disjoint; every input depends
only on the old variables. Rank means the dimension of the affine input
span. A nonzero span is called proper if it does not contain one.

\subsection{Affine linear algebra and ordinary restriction}

\begin{lemma}[Affine spans and coordinates]\label{lem:affine-coordinates}
\leangroup{\leanref{claims/AffineSystemLinearAlgebra.lean}{linear algebra}
  \leansep \leanref{claims/FiniteAffineSpan.lean}{affine spans}
  \leansep \leanref{claims/FiniteAffineCoordinates.lean}{coordinates}}
Let a finite affine system on \(\F^v\) have span \(S\).
Its common zero set is nonempty exactly when \(1\notin S\).
In that case it is an affine flat of codimension \(r=\dim S\), every
affine polynomial vanishing on it belongs to \(S\), and any basis
\(F_1,\ldots,F_r\) of \(S\) can be completed to invertible affine
coordinates. Both coordinate substitutions preserve ordinary degree.
If the zero set is empty, a constant linear combination of the inputs is one.
\end{lemma}
\begin{proof}
Write the inputs as \(g_i=c_i+L_i\). If a relation among their linear
parts had nonzero constant part, their span would contain one. Otherwise
all relations respect the right sides \(-c_i\); linear algebra then solves
\(L_i(x)=-c_i\). Equivalently, define the consistent functional on the
span of the \(L_i\) and extend it to the full dual space.
When \(1\notin S\), independence of the \(F_j\) implies independence
of their linear parts: a relation would be a constant in \(S\), hence
zero, and then all coefficients vanish. Complete those linear parts to a
basis and retain the constants of the \(F_j\). This gives an invertible
affine coordinate change. The zero flat has first \(r\) coordinates zero,
so an affine polynomial vanishing there is a linear combination of those
coordinates. Substitution and its inverse have degree at most one,
which proves degree preservation. Inconsistency is exactly the failed
relation already identified, normalized to give one.
\end{proof}

\begin{lemma}[Scalar cleanup and exact block degrees]\label{lem:cleanup}
\leangroup{\leanref{claims/FreshENSScalarCleanup.lean}{cleanup}
  \leansep \leanref{claims/FreshENSBlock.lean}{exact degrees}}
Over \(\F\), zero-span blocks, and unit-span blocks with \(h\ge1\), can
be removed by constant coefficient substitutions without increasing PC degree.
If a block has a genuine degree-one input, then
\(\deg P=2h\); each nonzero degree-one input has companion degree \(2h+1\).
In particular these degree equalities hold after discarding zero inputs in
a proper nonzero block. Without these hypotheses they are only upper bounds.
\end{lemma}
\begin{proof}
For zero span set all coefficients to zero. For unit span choose
\(\sum_i c_i g_i=1\), use these constants in the first coefficient row,
and use zero in the other rows. The product and companions become zero;
over \(\F\), every coefficient satisfies \(c_i^2=c_i\).
Apply the substitution lemma with degree multiplier one to the remaining
axioms. For the degree assertion, a genuine linear part contributes a
nonzero degree-two term to each factor, since its coefficient variables
are distinct. The polynomial ring is a domain, so degrees add in the
product and on multiplication by a degree-one input.
\end{proof}

\begin{lemma}[Ordinary restriction dimension and ideal membership]
\label{lem:restriction-ideal}
\leangroup{\leanref{claims/OrdinaryRestrictionDimension.lean}{dimension}
  \leansep \leanref{claims/AffineRestrictionIdeal.lean}{ideal membership}}
An affine substitution into \(d\) variables maps a space
of polynomials of degree at most \(k\) to a space of dimension at most
\(\binom{d+k}{k}\). Suppose the common zero flat of a finite affine
system \(g_i\) is nonempty. If the ordinary polynomial restriction of
\(f\), of degree at most \(k\), to that flat is zero, then for \(k\ge1\)
\begin{equation}\label{eq:literal-ideal}
 f=\sum_i a_i g_i,\qquad \deg a_i\le k-1.
\end{equation}
Here restriction means substitution into affine free coordinates, not
merely evaluation at field points.
\end{lemma}
\begin{proof}
Substitution does not increase degree. The target monomials of degree at
most \(k\), counted by \(\binom{d+k}{k}\), span the image.
For ideal membership use the coordinates of Lemma~\ref{lem:affine-coordinates}.
A polynomial whose restriction at \(F_1=\cdots=F_r=0\) is literally zero
has each monomial divisible by at least one \(F_j\). Assign each monomial
to one such factor and remove it. The resulting cofactors have degree at
most \(k-1\). Pull back the affine coordinate change and express each
\(F_j\) as a constant combination of the original inputs.
\end{proof}

\paragraph{Idea of the common kernel.}
The next lemma finds one row-linear \(f\) whose restriction to the zero flat
of every proper high-rank block is the zero polynomial. It is a dimension
count. Restriction to a flat of codimension \(r\) is a linear map from
\(\mathcal L_k\) into polynomials of degree at most \(k\) in \(v-r\) free
coordinates, a space of dimension \(\binom{v-r+k}{k}\). For
\(r>3\ell(k+1)\) this is smaller than \(\dim\mathcal L_k\ge\binom mk\ell^k\)
by a factor \(e^{-k^2/m}\). If \(M\) such maps are given and
\(Me^{-k^2/m}<1\), their joint kernel is nonzero.

\begin{lemma}[A common ordinary restriction kernel]\label{lem:kernel}
\leangroup{\leanref{claims/RestrictionKernelBounds.lean}{parameter bounds}
  \leansep \leanref{claims/CommonAffineRestrictionKernel.lean}{common kernel}}
Let \(m>0\), \(\ell\ge1\), \(h=3\ell\), \(M,k\ge1\), \(k\le m\), and
\[
 m\ln(4M)\le k^2.
\]
In a finite family of affine blocks in \(v=m\ell\) bits, suppose at most
\(M\) blocks are proper and have rank greater than \(h(k+1)\).
There is a nonzero \(f\in\mathcal L_k\) such that every block above
that rank threshold either has one in its input span or admits
\eqref{eq:literal-ideal}. In each proper high-rank block, the ordinary
restriction of \(f\) to its zero flat is zero.
\end{lemma}
\begin{proof}
Put \(r_*=3\ell(k+1)+1\). If \(r_*>v\), no proper high-rank block
exists and choose \(f=1\). Otherwise its restriction uses at most
\(d=v-r_*\) free coordinates. The checked dimension comparison uses
\[
 k!\binom{d+k}{k}\le(d+k)^k,\qquad
 k!\binom mk\ge(m-k+1)^k.
\]
Moreover \(d+k\le\ell(m-2k)\le\ell(m-k+1)(1-k/m)\).
The first inequality follows by inserting \(d=m\ell-3\ell(k+1)-1\)
and using \(\ell\ge1\); the second follows by expansion and \(k\le m\).
Thus
\[
 \frac{\binom{d+k}{k}}{\binom mk\ell^k}
 \le\left(1-\frac{k}{m}\right)^k\le e^{-k^2/m}.
\]
Consequently the sum of the restriction-image dimensions is at most
\[
 M e^{-k^2/m}\binom mk\ell^k
 \le\tfrac14\binom mk\ell^k<\dim\mathcal L_k.
\]
The joint restriction map has a nonzero kernel. Apply
Lemma~\ref{lem:restriction-ideal} to each proper high-rank block;
unit-span blocks need no restriction condition.
\end{proof}

\begin{corollary}[A convenient kernel parameter]\label{cor:kernel-parameter}
\leangroup{\leanref{claims/RestrictionKernelBounds.lean}{parameter bounds}
  \leansep \leanref{claims/CommonAffineRestrictionKernel.lean}{common kernel}}
For \(m,M\ge1\), the choice \(k=\lceil\sqrt{m\ln(4M)}\rceil\) is
positive and satisfies the square condition of Lemma~\ref{lem:kernel}.
Thus that lemma applies whenever additionally \(k\le m\).
\end{corollary}
\begin{proof}
The radicand is positive; square the inequality defining the ceiling.
\end{proof}

\subsection{One simultaneous weighted substitution}

The two substitutions of this subsection were described in
Section~\ref{sec:overview}: a low-rank block is replaced exactly by the
indicator of its zero flat, a high-rank block by \(1-f\), and the whole
refutation is replayed with weight \(f\). The degree accounting below shows
that every axiom image is derivable from the old base through \(kD+\deg f\),
after which Lemma~\ref{lem:substitution} does the rest.

\begin{lemma}[Low-rank packing without a retained core]\label{lem:packing}
\lean{claims/LowRankENS.lean}
For affine inputs over \(\F\) of span rank \(r\le h(k+1)\), with
\(h,k\ge0\), there are coefficient polynomials of degree at most \(k\)
substituting the block product to
\[
 Z=\prod_{j=1}^r(1-F_j),
\]
where \(F_j\) is a basis of the input span. This polynomial has degree
at most \(r\), is the Boolean indicator of the common zero set, and
\(g_iZ\in\Ideal_{r+1}(\Bool)\) for every input. No properness assumption
is needed.
\end{lemma}
\begin{proof}
Partition the basis into \(h\) ordered bins of size at most \(k+1\).
For each bin \(J\), use
\[
 1-\prod_{j\in J}(1-F_j)
   =\sum_{j\in J} F_j\prod_{t\in J,\ t<j}(1-F_t).
\]
Express each \(F_j\) as a constant combination of the inputs. Every
coefficient has at most \(k\) affine factors. Multiplying the resulting
bin products gives \(Z\). At every Boolean point it is one precisely
when the whole input span vanishes, so \(g_iZ\) vanishes everywhere and
has degree at most \(r+1\). Apply Boolean reduction. If \(h=0\), the
rank bound forces \(r=0\), all inputs are zero, and the empty product is one.
\end{proof}

\begin{lemma}[Weighted removal, including unit-span blocks]\label{lem:hybrid}
\lean{claims/AffineFamilyRemovalWithUnits.lean}
Let an old system \(\mathcal F\) in finitely many variables over \(\F\)
contain all old Boolean equations. Adjoin complete accuracy-\(h\) blocks
with finite affine input tuples and disjoint fresh coefficients, where
\(h,k\ge1\). Suppose this system has a degree-\(D\) ordinary-PC refutation,
\(D\ge2h+1\). Let \(f\) be any ordinary polynomial of degree \(s\le k\).
For every block of rank greater than \(h(k+1)\), suppose either its input
span contains one or it admits \eqref{eq:literal-ideal}. Then
\[
 f\in\Cons_{k(D+1)}(\mathcal F).
\]
No multilinearity or properness is required of \(f\) or of the blocks.
\end{lemma}
\begin{proof}
Make one simultaneous substitution fixing the old variables, with all
coefficient images of degree at most \(k\). For a low-rank block use
Lemma~\ref{lem:packing}. Its weighted companions have Boolean certificates
through \(s+r+1\le s+h(k+1)+1\le s+kD\).
For a high unit-span block, use the constant substitution making its
product zero. For every other high block use the cofactors of
\eqref{eq:literal-ideal} in the first row and zero elsewhere. Its product
becomes \(1-f\), and a weighted companion becomes
\[
 fg_i(1-f)=-g_i(f^2-f).
\]
Its Boolean certificate costs at most \(2s+1\le s+kD\).
A coefficient image \(\beta\) has weighted Boolean equation
\(f(\beta^2-\beta)\), certified through \(s+2k\le s+kD\).
An old axiom of degree \(e\le D\) has weighted image certified through
\(s+e\le s+kD\). These estimates initialize every source axiom used
by the refutation at the common ceiling \(kD+s\), irrespective of any
companion whose actual degree is below its general upper bound.
Lemma~\ref{lem:substitution}, with \(T=k\) and weight \(f\), replays the
whole primitive derivation. The final image of one is \(f\).
\end{proof}

\begin{theorem}[Finite-parameter affine exclusion]\label{thm:affine}
\lean{claims/AffineFamilyExclusion.lean}
Let \(m>0\), \(\ell\ge2\), \(n=2^\ell\), \(h=3\ell\), and let a finite
complete old-affine family be adjoined to \(\mathcal Q_{m,\ell}\).
Suppose \(M,k\ge1\), at most \(M\) blocks are proper and of rank greater
than \(h(k+1)\), and, with \(B=k(D+1)\),
\begin{equation}\label{eq:affine-room}
 m\ln(4M)\le k^2,\quad k\le m,\quad D\ge2h+1,\quad
 2B-1\le n,\quad4(k-1)<n.
\end{equation}
Then the augmented system has no ordinary-PC refutation through degree \(D\).
\end{theorem}
\begin{proof}
The common-kernel lemma supplies nonzero \(f\in\mathcal L_k\).
Weighted removal would derive \(f\) from the old base through \(B\),
contrary to Theorem~\ref{thm:row-separation}. Zero and unit spans are
already covered by removal, so no preliminary change of the family is needed.
\end{proof}

\begin{lemma}[Uniform eventual parameter room]\label{lem:parameters}
\lean{claims/AffineParameterBounds.lean}
Fix natural numbers \(a,A,d\). Put \(n=2^\ell\), \(t=\ell+1\),
\(q=\sqrt2\), and \(k=\lceil t q^\ell\rceil\).
For every sufficiently large \(\ell\), simultaneously
\[
 \ell\ge2,\quad 1\le k\le n+1,\quad
 (n+1)\ln\bigl(4(A2^{a\ell}+1)\bigr)\le k^2,
\]
\[
 4(k-1)<n,\qquad 2k\bigl(A(\ell+1)^d+1\bigr)-1\le n.
\]
The threshold depends only on \(a,A,d\).
\end{lemma}
\begin{proof}
The ceiling gives \(tq^\ell\le k\le2tq^\ell\).
Set \(C=\ln(4(A+1))\). Since \(\ln2\le1\),
\(\ln(4(A2^{a\ell}+1))\le C+a\ell\).
For sufficiently large \(\ell\), \(2(C+a\ell)\le t^2\).
Together with \(n+1\le2n\) and \(q^{2\ell}=n\), this proves the
square condition. For every fixed natural \(e\),
\(t^e/q^\ell\longrightarrow0\): its successive ratio tends to
\(1/q<1\), so its tail is dominated by a geometric sequence.
Applying this to \(e=d+1\) gives eventually
\[
 4k(At^d+1)\le8(A+1)t^{d+1}q^\ell<n.
\]
This implies the remaining degree and packing bounds, and also \(k\le n+1\).
\end{proof}

\begin{corollary}[Polynomial inventory and polylogarithmic degree]
\label{cor:polylog-exclusion}\lean{claims/AffineFamilyExclusion.lean}
Fix natural \(a,A,d\). For all sufficiently large \(\ell\), every finite
complete accuracy-\(3\ell\) affine family over \(\mathcal Q_{2^\ell+1,\ell}\)
with at most \(A2^{a\ell}+1\) proper blocks has no PC refutation of degree
\(D\le A(\ell+1)^d\).
\end{corollary}
\begin{proof}
Raise the proposed degree ceiling to \(D'=\max\{D,6\ell+1\}\).
It is at most \((A+7)(\ell+1)^{d+1}\). Apply
Lemma~\ref{lem:parameters} with constants \(a,A+7,d+1\), and use its
\(k\) and the upper inventory bound in Theorem~\ref{thm:affine}.
All required inequalities hold uniformly over the actual families.
\end{proof}
